Benchmark contamination -- the presence of evaluation data in a large language model's (LLM) training corpus -- inflates reported scores on static test suites in ways that are difficult to detect after the fact. This work evaluates whether a dynamic, template-based benchmark for LLM unit-test generation resists this inflation more effectively than a state-of-the-art static alternative. Sixty parameterised templates spanning four cyclomatic-complexity tiers generate fresh, seeded problem instances at evaluation time, scored by a seven-operator mutation-testing suite rather than fixed reference tests. Nine models across four providers are evaluated against the dynamic benchmark and a decontaminated static baseline (ULT), and a small open-source model is fine-tuned directly on the static baseline's ground truth to causally probe contamination sensitivity. The fine-tuned model gains \HTwoSeenGain{} percentage points on the memorised static split against a \HTwoLOneLThreeGain{}-point generic-improvement control, while its dynamic-benchmark score moves by exactly zero -- the clearest evidence obtained that the dynamic benchmark rejects contamination the static one absorbs. Six pre-registered hypotheses are evaluated under Holm-Bonferroni correction; results are mixed on statistical significance but consistent in direction, and the dynamic benchmark tracks an external code-generation proxy at least as well as the static benchmark does. The findings support dynamic, template-based generation as a practical contamination-resistant evaluation strategy for LLM test-generation research, while identifying where its statistical power currently falls short of a pre-registered bar.
